% GENERATED FILE. Edit canonical lessons, then run npm run generate:books.
% canonical-source-hash: fnv1a64:2606e98487e2c8e6
% canonical-lessons: ML-C26-raavile

\chapter{Morning}
\label{ch:ml-morning}

This chapter is generated from the canonical micro-lessons used by Language Ladder.
Each section stays independently resumable and preserves the authored prerequisite order.

\section[\ml{രാവിലെ}]{\ml{രാവിലെ} (rāvile) — "morning," built from "night"}

\label{lesson:ML-C26-raavile}



\begin{quote}
\textbf{Warm-up.} {[}PAUSE 2s{]} Kannada's morning word means "the shining." Telugu's means "rise." Malayalam's most common morning word means something stranger on its face: night.
\end{quote}

\subsection*{You'll want to know}
\textbf{\ml{രാവിലെ}} (\textbf{rāvile}) — "\textbf{morning}" — is, per Wiktionary, a compound: \textbf{\ml{രാവ്}} (\emph{rāvŭ}) plus the suffixes \textbf{-il} and \textbf{-e}. Here's the honest surprise: \textbf{\ml{രാവ്}} shares its root — not its exact form — with the word already met in Chapter 24, \textbf{\ml{ഇരവ്}} (\emph{iravŭ}), there a \textbf{literary-register synonym for "night" itself}. Wiktionary confirms both as cognates, both tracing to \textbf{Proto-Dravidian} \emph{\textbf{*cira}} ("\textbf{darkness}") — a cognate doublet within Malayalam, not literally one word in two spellings. So Malayalam's everyday word for "morning" is, transparently, built from \textbf{night} plus locative-ish suffixes — not from light, dawn, or rising, the strategies Kannada's \textbf{\ka{ಬೆಳಗ್ಗೆ}} ("the shining") and Telugu's \textbf{\te{ఉదయం}} (Sanskrit "rise") both reach for instead.

\begin{culture}
No source spells out \textbf{why} "at/in the night" came to mean "morning" specifically — it's a plausible semantic path (perhaps "the night just ending," the pre-dawn hours) rather than one directly documented step by step. Treat the \textbf{morphological composition} as confirmed (Wiktionary is explicit about the rāvŭ+il+e breakdown); treat the \textbf{semantic mechanism} as a reasonable but unconfirmed inference, not a settled fact.
\end{culture}

\begin{culture}
Malayalam keeps two other words for "morning" too. \textbf{\ml{പ്രഭാതം}} (\emph{prabhātaṁ}) is a Sanskrit \textbf{tatsama} sharing its root — \textbf{prabhāta}, "to shine, become bright" — with Telugu's own \textbf{\te{సుప్రభాతం}} (\emph{suprabhātam}, TE-C29, this very arc's Telugu morning-greeting word) — a genuine cross-language callback within this project. \textbf{\ml{പുലർച്ച}} (\emph{pulaṟcca}) is native Dravidian, but be honest about its etymology: one database entry links it to a root meaning "to subsist, live," which may or may not be the same root behind its "dawn" sense — treat that specific connection as unconfirmed, not settled.
\end{culture}

\subsection*{Guided Practice}
{[}PAUSE 1s{]}

\begin{itemize}
\raggedright
  \item {[}YOU SAY: "rāvile" — "morning," built from rāvŭ ("night") + locative suffixes{]}
  \item {[}YOU SAY: the honest surprise — a cognate doublet of the night-word already met as iravŭ (Chapter 24), same root, now doing morning's job{]}
  \item {[}YOU SAY: prabhātaṁ — same Sanskrit root as Telugu's suprabhātam{]}
\end{itemize}

\begin{tcolorbox}[breakable,colback=teal!4,colframe=teal!35!black,title={Wrap-up recall}]
{[}PAUSE 3s{]} What is \ml{രാവിലെ} literally built from? (\textbf{\ml{രാവ്}} — "night" — plus locative suffixes -il/-e; a cognate doublet, same root, as \ml{ഇരവ്} met in Chapter 24 — not literally the same word.) Is the exact reason "night" came to mean "morning" fully documented? (\textbf{No} — the compound itself is confirmed, but the semantic shift is a plausible, unconfirmed inference.) What root does \ml{പ്രഭാതം} share with a word from Telugu's arc? (\textbf{prabhāta} — the same root behind Telugu's \te{సుప్రభాతం}, TE-C29.)
\end{tcolorbox}
